\documentclass[10pt,a4paper]{amsart}
\usepackage[margin=19mm]{geometry}
\usepackage{amsmath,amssymb,mathtools}
\usepackage[scr=rsfs,cal=euler]{mathalfa}
\usepackage{xfrac}
\usepackage[unicode,hidelinks,bookmarks=false]{hyperref}
\newcommand{\ie}{i.e.\ }
\newcommand{\cf}{cf.\ }
\newcommand{\resp}{respectively\ }
\newcommand{\Subsection}{\subsection}
\providecommand{\nobreakdash}{\nobreak\hbox{-}}
\setlength{\emergencystretch}{2em}
\multlinegap0pt
\usepackage{bm}
\usepackage{bbm}
\usepackage{dsfont}
\usepackage{xspace}
\usepackage{cancel}
\usepackage{mathtools}
\usepackage{amsthm}
\makeatletter
\@ifundefined{theorem}{\newtheorem{theorem}{Theorem}}{}
\@ifundefined{lemma}{\newtheorem{lemma}[theorem]{Lemma}}{}
\@ifundefined{proposition}{\newtheorem{proposition}[theorem]{Proposition}}{}
\makeatother
\def\Ker{\mathrm{Ker \,}}
\def\K{\mathcal{K}}
\newcommand{\Cinfty}{\mathscr{C}^\infty}
\newcommand{\Lie}{\mathscr{L}}
\newcommand{\X}{\mathfrak{X}}
\newcommand{\inn}{i}
\renewcommand{\d}{\mathrm{d}}
\title[Pairs of differential forms: a framework for precontact geom]{Pairs of differential forms: a framework for precontact geometry}
\author{Xavier Gr\`acia, \`Angel Mart\'inez-Mu\~noz, Xavier Rivas}
\date{Source: arXiv:2602.04882v1 (CC BY 4.0)}
\begin{document}
\maketitle

\noindent\emph{Source and licence.} Pairs of differential forms: a framework for precontact geometry. Version arXiv:2602.04882v1 (CC BY 4.0), distributed under the Creative Commons Attribution 4.0 International licence, as recorded in the arXiv licence metadata for this exact version and shown on the abstract page https://arxiv.org/abs/2602.04882v1. Full rights notice: licenses/arxiv-2602.04882-CC-BY-4.0.txt.\par\smallskip

\noindent\textit{Editorial scope.} Counted unit: the Proposition whose source label is \texttt{change\_of\_class\_parity} in the article (source0.tex lines 2149--2151, source label \texttt{change\_of\_class\_parity}), delivered together with the article's own complete original proof of it (source0.tex lines 2152--2190, one \texttt{proof} environment of 39 lines). No mathematics is written, repaired or re-derived: statement and proof are the article's own text line for line. Required context retained uncounted: conditions\_odd\_class (lines 980--988); eq:even\_conditions (lines 1061--1063); prop:liouville\_structure (lines 1177--1179); lemma\_formula\_ef\_times\_form (lines 2085--2090). Every definition, display and label of those lines is retained unchanged. Omitted from this package and not redistributed: 1--979, 989--1060, 1064--1176, 1180--2084, 2091--2148, 2191--2436. Nothing inside a retained excerpt is omitted or altered: every retained body is delivered exactly as the article's lines are written, so no omission receipt of any kind accompanies this package. Citations: the retained excerpts carry no citation command at all, so no bibliography entry of the article is transcribed and no reference is redistributed. Reference adaptation: none was needed and none was made; every reference inside the delivered unit resolves inside it, so no unresolved reference or citation is produced. Printed slips kept literal: no printed word, symbol, hypothesis or number of the counted statement or of its proof was altered. Layout and class: the article's own class is \texttt{article} with packages including inputenc, colortbl, graphicx, srcltx, amsmath, amssymb, bm, amsthm, xy, mathrsfs, mathtools, tikz-cd, xcolor, enumitem; the wrapper is the lane's pinned \texttt{amsart} editorial wrapper at 10pt on A4, which restates the theorem-like environments the retained text needs and adds the article's own macro definitions that the retained text needs (\texttt{\textbackslash Cinfty}, \texttt{\textbackslash K}, \texttt{\textbackslash Ker}, \texttt{\textbackslash Lie}, \texttt{\textbackslash X}, \texttt{\textbackslash d}, \texttt{\textbackslash inn}), with extra wrapper packages (bm, bbm, dsfont, xspace, cancel, mathtools). The delivered numbering is the unit's own. Assets: no retained span contains a figure environment, a graphics-inclusion command, a PDF-page inclusion, a table or a TikZ picture; no image file is copied into this package, the package declares no asset dependency and no image file is redistributed with it. Licence: version arXiv:2602.04882v1 of this article is distributed under the Creative Commons Attribution 4.0 International licence, as recorded in the arXiv licence metadata for this exact version and shown on the abstract page https://arxiv.org/abs/2602.04882v1. A full rights notice is delivered at licenses/arxiv-2602.04882-CC-BY-4.0.txt. Characters: every retained excerpt is pure ASCII, so the delivered engine, XeLaTeX with its default Latin Modern text font, reports no missing character.
\begin{proposition}\label{conditions_odd_class}
   Let $(\tau,\omega)$ be a pair of a differential 1-form and a 2-form.  
Then $(\tau,\omega)$ is a doublet of class $2r+1$ if, and only if,
    \begin{equation}\label{eq:odd_conditions}
            \omega ^{r+1} = 0\,,
            \qquad
            \tau \wedge \omega^r \neq 0\,.
    \end{equation}
\end{proposition}

    \begin{equation}\label{eq:even_conditions}
\omega^r \neq 0\,, \qquad \omega ^{r+1}=0\,, \qquad \tau\wedge \omega^r=0\,.
    \end{equation}

\begin{proposition}\label{prop:liouville_structure}
    Let $(\tau,\omega)$ be a doublet of even class. Then, a vector field $X$ is a Liouville vector field if and only if $\widehat B(X) = \tau$. Hence, all Liouville vector fields are given by $\Delta=\Delta_0+\Gamma$, where $\Delta_0$ is a particular Liouville vector field and $\Gamma\in\K_{(\tau,\omega)}=\Ker\omega$.
\end{proposition}

\begin{lemma} 
\label{lemma_formula_ef_times_form}
Let $\tau\in\Omega^1(M)$ be a $1$-form on $M$ and $f\in \Cinfty(M)$ a function. 
Then, for every $n \geq 1$, 
$$ (e^f(\d f\wedge \tau +\omega))^n = e^{nf}(n\, \d f \wedge\tau\wedge\omega^{n-1}+\omega^n)\,.$$
\end{lemma}

\begin{proposition}\label{change_of_class_parity}
    Let $(\tau,\omega)$ be a doublet of class $2r+2$, with $\tau$ nowhere-vanishing, and $f\in\Cinfty(M)$ a function. Then, $(e^f\tau, e^f(\d f\wedge \tau+\omega))$ is a doublet of class $2r+1$ if, and only if, for all Liouville vector fields, we have $$\Lie _{\Delta} f=-1\,.$$
\end{proposition}

\begin{proof}
   Let us first assume that 
   $(e^f\tau, e^f(\d f\wedge \tau+\omega))$ 
   is a doublet of class $2r+1$. We will denote $\widetilde \omega \coloneqq \d f \wedge \tau+\omega$. As the doublet $(\tau,\omega)$ has even class, there must exist Liouville vector fields $\Delta \in \X(M)$.
    If we contract $(e^f\widetilde \omega)^{r+1}$ by any Liouville vector field $\Delta$ of $(\tau,\omega)$, we obtain
\begin{align}\label{Liouville_function_class}
    \begin{aligned}
        \inn{\Delta}(e^f\widetilde \omega )^{r+1} &=  e^{(r+1)f}((r+1)\, \inn{\Delta}(\d f \wedge\tau\wedge\omega^{r})+\inn{\Delta}\omega^{r+1})
        \\
        &=e^{(r+1)f}(((r+1)\Lie_{\Delta}f)\,\tau\wedge\omega^{r} + (r+1) \,\tau\wedge\omega^{r})
        \\
        &=(e^{(r+1)f}(r+1)(\Lie_{\Delta} f +1 ))\,\tau\wedge\omega^{r}\,.
        \end{aligned}
    \end{align}
    where we have used Lemma \ref{lemma_formula_ef_times_form} in the first equality.
Now, as the class of $(\tau,\omega)$ is $2r+2$, and the class of $(e^f \tau, e^f \widetilde \omega )$ is $2r+1$, by Propositions~\ref{conditions_odd_class} and~\ref{eq:even_conditions}, we have
$(e^f\widetilde \omega)^{r+1}=0$ and $ \qquad\tau\wedge\omega^r\neq 0$. 
Therefore, it follows from equation \eqref{Liouville_function_class} that 
$ \Lie_{\Delta} f +1=0 $,
for all Liouville vector fields $\Delta$ of the doublet $(\tau,\omega)$, as we wanted to see.

To see the converse first recall that all the Liouville vector fields $\Delta\in\X(M)$ of $(\tau,\omega)$
can be written as $\Delta = \Delta_0+\Gamma$, where $\Delta_0\in\X(M)$ is a particular Liouville and $\Gamma\in \Ker\omega$ (see Proposition~\ref{prop:liouville_structure}). Thus, if the assumption is satisfied for every Liouville vector field, one has that
$$\inn{\Delta_0}\d f+1=0\qquad \text{and} \qquad \inn{\Delta_0+\Gamma}\d f +1=\inn{\Delta_0}\d f+\inn{\Gamma}\d f+1=0\,,$$
which implies that 
$\inn{\Gamma}\d f = 0$,
for all $\Gamma \in \Ker\omega$. 
Using this and $\Ker\omega \subset \Ker \tau$ (due to the class being even), it is easy to check that
$\inn{\Gamma} (e^f\widetilde \omega) = 0$,
and also that
$ \inn{\Delta}(e^f\widetilde \omega) = e^f((\inn{\Delta} \d f)\tau + \tau) = 0$.
Thus, we have
$ \mathcal{X}_{(\tau,\omega)} =\K_{(\tau,\omega)} \oplus \left\langle \Delta_0 \right\rangle  \subseteq \Ker (e^f \tau )\cap \Ker (e^f\widetilde \omega)$,
but as the class of the modified pair cannot increase or decrease by more than 1, we necessarily have
$$\mathcal{X}_{(\tau,\omega)}=\K_{(\tau,\omega)} \oplus \left\langle \Delta_0 \right\rangle  = \Ker (e^f \tau )\cap \Ker (e^f\widetilde \omega)\,,$$
which is the characteristic distribution of the conformally transformed pair
$(e^f\tau,e^f\widetilde \omega)$;
therefore its class is $2r+1$.
\end{proof}

\end{document}
